% Generated draft: requires mathematical and source review.
\documentclass[11pt]{article}
\usepackage[margin=25mm]{geometry}
\usepackage{fontspec}
\setmainfont{Latin Modern Roman}
\usepackage{amsmath,amssymb,mathtools,unicode-math}
\setmathfont{Latin Modern Math}
\usepackage{enumitem,needspace,xurl,hyperref}
\setlength{\parindent}{0pt}
\setlength{\parskip}{5pt}
\setlength{\emergencystretch}{3em}
\newcommand{\N}{\mathbb{N}}
\newcommand{\Z}{\mathbb{Z}}
\newcommand{\Q}{\mathbb{Q}}
\newcommand{\R}{\mathbb{R}}
\newcommand{\C}{\mathbb{C}}
\newcommand{\F}{\mathbb{F}}
\newcommand{\A}{\mathbb{A}}
\newcommand{\PP}{\mathbb{P}}
\newcommand{\E}{\mathbb{E}}
\newcommand{\eps}{\varepsilon}
\newcommand{\dd}{\,\mathrm{d}}
\begin{document}
\section*{UnsolvedMath ID: 30006516}
\textbf{Source problem code:} catalogue entry
\subsection{Definitions and mathematical statement}
In the usual countably infinite setting, let \(A\) be a first-order structure on a countably infinite set \(\Omega\), in a countable language. Say that \(A\) is \(\omega\)-categorical if its complete theory has exactly one countably infinite model up to isomorphism. These countability conventions are an explicit interpretation of the terminology; the supplied excerpt does not state them separately.

For a finite set \(F\subseteq\Omega\), let \(\operatorname{acl}_A(F)\) consist of those elements belonging to a finite set definable in \(A\) with parameters from \(F\). The condition of \emph{no algebraicity} is
\[
\operatorname{acl}_A(F)=F\qquad\text{for every finite }F\subseteq\Omega,
\]
including \(F=\varnothing\). Let \(\operatorname{Aut}(A)\) be the group of automorphisms of \(A\). Let \(\operatorname{EEmb}(A)\) be the monoid, under composition, of elementary embeddings \(e:A\to A\): for every first-order formula \(\varphi(\bar x)\) and every suitable finite tuple \(\bar a\) from \(A\),
\[
A\models\varphi(\bar a)\quad\Longleftrightarrow\quad A\models\varphi(e(\bar a)).
\]
Both act naturally on \(\Omega\).

To specify reconstruction precisely, a class \(\mathcal C\) of comparison structures must be fixed. Write \(M(A)=\operatorname{Aut}(A)\) or \(M(A)=\operatorname{EEmb}(A)\), treating the two choices separately. For \(B\in\mathcal C\) with underlying set \(\Lambda\), a bijection \(h:\Omega\to\Lambda\) induces conjugation
\[
c_h(f)=h\circ f\circ h^{-1}.
\]
One precise existence formulation of reconstruction is
\[
\forall B\in\mathcal C,\qquad
M(A)\cong M(B)\ \Longrightarrow\
\exists h:\Omega\xrightarrow{\sim}\Lambda\quad
hM(A)h^{-1}=M(B).
\]
A stronger formulation requires every abstract group isomorphism, respectively monoid isomorphism, \(\Phi:M(A)\to M(B)\) to be induced by a bijection:
\[
\forall B\in\mathcal C\ \forall\Phi:M(A)\xrightarrow{\sim}M(B)\ \exists h:\Omega\xrightarrow{\sim}\Lambda\ \forall f\in M(A),\qquad
\Phi(f)=h\circ f\circ h^{-1}.
\]
These formulations distinguish reconstruction of the action from realization of each specified isomorphism.

\textbf{Source question.} Does every \(\omega\)-categorical structure \(A\) without algebraicity have automatic action reconstruction for its automorphism group or its elementary-embedding monoid? The supplied excerpt asks for reconstruction up to conjugation by a bijection, but does not specify \(\mathcal C\), define the word ``automatic,'' or resolve which of the displayed quantifier conventions is intended. Thus the displayed statements describe precise possible formulations, rather than a uniquely determined formalization of the source.

The source also describes the intended outcome as recovery of \(A\) up to bi-definability. Under the parameter-free convention, this means that, after transporting a comparison structure along a bijection, each structure's basic relations, function graphs, and named constants are definable without parameters in the other. The supplied excerpt does not state the hypotheses or conventions under which this conclusion is intended to follow from action reconstruction.

The supplied dated curation labels the problem partially solved: it reports a solution for elementary-embedding monoids under the stated no-algebraicity setting, while retaining the general automorphism-group question as open. This is a supplied literature assessment, not an independently verified theorem or status determination.
\subsection{Short English statement}
The question asks whether the abstract automorphism group or elementary-embedding monoid determines the underlying action of an \(\omega\)-categorical structure without algebraicity, up to relabelling its elements, with bi-definability as the stated intended consequence. The comparison class and the exact reconstruction convention remain unspecified. Supplied curation reports the monoid case solved and the general group case open, subject to independent verification.
\subsection{Sources}
\begin{itemize}
\item Supplied local UnsolvedMath registry record, permanent ID \(30006516\), titled \emph{Automatic Action Reconstruction for Countably Categorical Structures}; exact source problem code \texttt{catalogue entry}. The record belongs to the Oberwolfach Reports Open Problems collection. No individual proposer is supplied.
\item Supplied primary-source citation: \emph{Homogeneous Structures: Model Theory meets Universal Algebra} (2026), Oberwolfach Reports, DOI \texttt{10.4171/owr/2025/58}, pp.\ 3063--3124; the original excerpt identifies p.\ 3093. \url{https://doi.org/10.4171/owr/2025/58}. The registry separately records proposed year 2025 and identifies the report as 2025/58; these supplied year descriptions have not been reconciled.
\item Supplied literature source: de la Nuez Gonzalez et al., \emph{Minimal and intrinsic topologies on monoids}. \url{https://arxiv.org/abs/2603.28419}.
\item Supplied literature source: \emph{Oberwolfach Report 2025/58}. \url{https://doi.org/10.4171/owr/2025/58}.
\item Literature assessment date, copied exactly: \texttt{2026-08-22}. This is supplied dated curation derived from the supplied OWR curation, not independently verified. It labels the problem partially solved, reports automatic action reconstruction for elementary-embedding monoids under the stated no-algebraicity setting, and describes the corresponding general automorphism-group problem as open. The record expressly requests independent verification of time-sensitive status and cited claims.
\item The registry labels the statement \texttt{corrected\_verified} and reports that the official report was checked, the question isolated from surrounding exposition, and its notation repaired. These are supplied provenance claims; this draft does not independently confirm them.
\end{itemize}
\end{document}
